\centerline{\bf \Huge Разнобой по процессам}

\begin{flushright}
        \scriptsize{}
\end{flushright}

\problem{1}
Неуловимый Джо со 100 долларами в кармане пришёл играть в рулетку. Он никогда не проигрывает на рулетке больше пяти раз подряд и никогда не ставит больше 100 долларов. При проигрыше Джо ничего не получает, при выигрыше - получает удвоенную ставку. Как ему выиграть 1000 000?

\problem{2}
Олежка любит квас. У него есть две одинаковые кружки, одна полная кваса, вторая пустая. Олежка хочет перелить часть кваса в пустую кружку и угостить Эльзятку. После каждого переливания он решает, что перелил чуть больше чем хотел, и начинает переливать квас обратно. Первым действием Олежка перелил половину кваса из первой кружки во вторую. Вторым действием он перелил треть кваса, бывшего во второй кружке, в первую. Третьим действием он перелил четверть кваса, бывшего в первой кружке, во вторую и т.д. Сколько кваса будет в каждой из кружек после 2023 переливаний?

\problem{3}
Эльзятка нарисовала на доске $N$ точек (никакие три не лежат на одной прямой) и соединила некоторые пары из них отрезками. Если какие-то два отрезка пересекаются, то Олежка может их стереть и нарисовать другие два отрезка с теми же концами. Докажите, что после нескольких таких операций пересекающихся отрезков не останется.

\problem{4}
Правильный треугольник со стороной 3 разбит на девять правильных треугольничков. В этих треугольничках изначально записаны нули. За один ход можно выбрать два числа, находящиеся в соседних по стороне клетках, и либо прибавить к обоим по единице, либо вычесть из обоих по единице. Петя хочет сделать несколько ходов так, чтобы после этого в клетках оказались записаны в некотором порядке последовательные натуральные числа $n, n+1, \ldots, n+8$. При каких $n$ он сможет это сделать?

\problem{5}
Несколько доминошек уложены в виде прямоугольника. Если две доминошки образуют квадрат, то их можно повернуть на $90^{\circ}$. Докажите, что можно все доминошки сориентировать одинаково.